% 19b-outside.tex: последнее кольцо. Добавлено 25.09.2026, после того как холодное прочтение нашло
% кульминацию документа, первые программы извне, рассказанной рано (раздел 6) и погребённой под оговорками.
\section{Последнее кольцо: чужие программы}
\label{sec:outside}

\analogy{Каждую стену, башню и книгу в этом городе построили те, кто в нём живёт, и проверили тоже
они. Однажды пришли путники, не строившие ничего из этого. Они несли собственные печати, говорили на
общем языке дороги и ничем не были обязаны городу. То, что они сделали у ворот, есть единственный
отчёт, который город написал не о себе сам.}

Каждое кольцо этой карты существует потому, что кольцо внутри него может быть верным и всё же
неспособным доказать то, что обществу необходимо знать. Последнее кольцо есть то, которого никакая
внутренняя работа дать не может: программы, написанные кем-то другим, не разделяющие авторской модели
того, что содержит правильный артефакт, и встречающие Полярис у одной из его дверей. Дверь есть
верификатор ОИД4ВП из раздела~\ref{sec:verifiers}, выбранный потому, что он говорит на протоколе, на
котором в действительности говорит сторонний кошелёк.

\figfit{0.56}{outside-witnesses}{Три вещи за пределами репозитория, которые его задействовали, и два
контроля, без которых первая из них ничего бы не значила. Кошелёк, написанный здесь никем,
предъявил удостоверение и был принят; тот же кошелёк, когда верификатор доверял другому ключу
эмитента, получил отказ; и повторное предъявление против уже отвеченного запроса получило отказ.
Размещённый пакет Фонда сам засчитал семь отказов и оставил четыре принятия на человека-рецензента;
более поздний прогон того же плана сертифицирован. Опубликованный набор векторов сходится с обоими
свидетелями подписи на примитиве.}{fig:outside-witnesses}

\subsection{Первый чужак: служба соответствия}

Фонд ОпенАйДи ведёт размещённый пакет проверок соответствия с планом тестов ровно для этой роли, в
котором пакет сам играет кошелёк и сам подписывает удостоверение. Он прогнал все одиннадцать модулей
плана против верификатора через открытый интернет, забирая объект запроса и отправляя ответ снаружи.
Семь отрицательных модулей, где пакет присылает предъявление, испорченное одним определённым способом,
засчитывает сама служба по отказу верификатора, и это ближайшее подобие внешнего противника, какое
этому проекту предлагали. Четыре положительных модуля заканчиваются в состоянии, которое пакет
оставляет человеку, чтобы тот засвидетельствовал успешную проверку. Более поздний прогон того же плана
подан на сертификацию, и Фонд числит \code{полярис-оид4вп 1.0.0rc7} как сертифицированный ОпенАйДи
по профилю верификатора ОИД4ВП 1.0 + ХАИП 1.0. Каждый прогон несёт два отрицательных контроля, по
одному в каждую сторону: верификатор, исправленный так, чтобы отказывать во всём, замечают
положительные модули, а исправленный так, чтобы принимать всё, замечают все семь отрицательных. \sWit

\subsection{Второй чужак: кошелёк}

Затем пришёл кошелёк. Неизменённый интерфейс кошелька walt.id, опубликованный образ с закреплённым
дайджестом, забрал подписанный объект запроса, проверил его против зарегистрированного якоря доверия,
сопоставил своё удостоверение с запросом, подписал токен привязки ключом, который сам создал и которого
этот репозиторий никогда не держал, зашифровал ответ и отправил его. Верификатор проверил всю цепочку
и принял его.

Сперва он отказал Полярису, и был прав.

Собственный генератор ключей верификатора произвёл сертификат подписи запросов вовсе без расширения
назначения ключа, так что сертификат, вся работа которого есть подпись объектов запроса, не был
отмечен как пригодный для подписи. Каждый наблюдатель внутри репозитория прошёл мимо: модули
соответствия, тесты пакета, мутационное учение над каждым отказом, каждая проверка слоя инвариантов.
Каждый из них наследует авторскую модель того, что такой сертификат обязан содержать. Первый
наблюдатель, её не унаследовавший, нашёл пробел сразу. Сертификат исправлен, с тестом, падающим на
листе без этого расширения, а обмен несёт те же два контроля, что и прогон соответствия, потому что
верификатор, принимающий всё, печатает ту же строку успеха. Любой может повторить обмен с чистой
машины примерно за десять минут, не клонируя репозиторий. \sWit

Это довод всей карты в одном событии. Кольца принадлежат автору; последнее кольцо принадлежит всем
остальным, и только оно способно увидеть то, что остальные разделяют.

\subsection{Чего последнее кольцо пока не устанавливает}

Сказано один раз и полностью. Сертификация есть самосертификация, которую Фонд рассмотрел и опубликовал,
для одного пакета, одной версии и одной роли: это не одобрение, не аудит и не утверждение об остальном
Полярис. Обмен с кошельком покрывает один кошелёк, один формат удостоверения, один путь предъявления и
один классический алгоритм подписи, как и закрепляет профиль; он ничего не говорит ни о другом
кошельке, ни о мобильных документах, ни о постквантовом пути, ни о совместимости вообще. Собственный
протокол, форматы \code{полярис-*/1}, на которых построена остальная карта, никто извне по одним
документам не реализовал. Ни один оператор, кроме автора, систему не запускал, и ни одна независимая
сторона её не проверяла. Каждое из этого есть строка в сводной таблице проекта, и каждая строка пуста,
пока её не заполнит названная сторона (раздел~\ref{sec:limitations}). \sExt
